% Bagian: first-order-logic
% Bab: proof-systems
% Subbagian: natural-deduction

\documentclass[../../../include/open-logic-section]{subfiles}

\begin{document}

\iftag{FOL}
      {\olfileid[id]{fol}{prf}{ntd}}
      {\olfileid[id]{pl}{prf}{ntd}}

\olsection{Deduksi Alami}

Deduksi alami adalah sistem !!a{derivation} yang dimaksudkan untuk
mencerminkan penalaran yang sesungguhnya (terutama jenis penalaran
sistematis yang digunakan oleh matematikawan). Penalaran yang
sesungguhnya berlangsung mengikuti sejumlah pola ``alami.'' Sebagai
contoh, pembuktian per kasus memungkinkan kita menetapkan suatu
kesimpulan berdasarkan premis disjungtif dengan menetapkan bahwa
kesimpulan itu mengikuti dari masing-masing disjung. Pembuktian tak
langsung memungkinkan kita menetapkan suatu kesimpulan dengan
menunjukkan bahwa negasinya mengarah pada kontradiksi. Pembuktian
bersyarat menetapkan klaim bersyarat ``jika \dots maka \dots'' dengan
menunjukkan bahwa konsekuen mengikuti dari anteseden. Deduksi alami
merupakan formalisasi dari beberapa inferensi alami ini. Setiap
penghubung logis dan kuantor memiliki dua aturan, yaitu aturan
introduksi dan aturan eliminasi, dan masing-masing bersesuaian dengan
salah satu pola inferensi alami tersebut. Sebagai contoh,
$\Intro{\lif}$ bersesuaian dengan pembuktian bersyarat, dan
$\Elim{\lor}$ dengan pembuktian per kasus. Aturan yang sangat sederhana
ialah $\Elim{\land}$, yang memungkinkan inferensi dari $!A \land !B$
ke~$!A$ (atau $!B$).

Salah satu ciri yang membedakan deduksi alami dari sistem
!!{derivation} lain ialah penggunaan asumsi. !!^a{derivation} dalam
deduksi alami merupakan pohon !!{formula}s. Sebuah !!{formula}
tunggal berada di akar pohon !!{formula}s tersebut, sedangkan
``daun-daun'' pohon itu adalah !!{formula}s yang darinya kesimpulan
diturunkan. Dalam deduksi alami, beberapa !!{formula}s daun berperan
di dalam !!{derivation}, tetapi sudah ``habis digunakan'' ketika
!!{derivation} mencapai kesimpulan. Hal ini bersesuaian dengan praktik
dalam penalaran yang sesungguhnya untuk memperkenalkan hipotesis yang
hanya berlaku dalam waktu singkat. Sebagai contoh, dalam pembuktian
per kasus kita mengasumsikan kebenaran masing-masing disjung; dalam
pembuktian bersyarat kita mengasumsikan kebenaran anteseden; dalam
pembuktian tak langsung kita mengasumsikan kebenaran negasi
kesimpulan. Cara memperkenalkan asumsi hipotetis, lalu melepaskannya
demi menetapkan suatu langkah perantara, merupakan ciri khas deduksi
alami. Formula pada daun suatu !!{derivation} deduksi alami
disebut asumsi, dan beberapa aturan inferensi dapat
``!!{discharge}'' asumsi-asumsi tersebut. Sebagai contoh, jika kita
memiliki !!a{derivation} untuk~$!B$ dari beberapa asumsi yang
mencakup~$!A$, aturan $\Intro{\lif}$ memungkinkan kita
menyimpulkan~$!A \lif !B$ dan melepaskan asumsi-asumsi berbentuk~$!A$
yang dipilih serta diberi label untuk inferensi tersebut.
(Untuk mencatat asumsi mana yang dilepaskan pada inferensi mana, kita
memberi label berupa angka pada inferensi dan asumsi yang
dilepaskannya.) Asumsi yang !!{undischarged} pada akhir
!!{derivation} secara bersama-sama cukup bagi kebenaran kesimpulan;
dengan demikian, !!a{derivation} menetapkan bahwa asumsi-asumsi yang
!!{undischarged} dalam derivasi itu mengakibatkan kesimpulannya.

Relasi $\Gamma \Proves !A$ berdasarkan deduksi alami berlaku jika dan
hanya jika terdapat !!a{derivation} dengan $!A$~sebagai !!{sentence}
terakhir dalam pohon, dan setiap daun yang !!{undischarged} berada
dalam~$\Gamma$. $!A$~merupakan teorema dalam deduksi alami jika dan
hanya jika terdapat !!a{derivation} dengan $!A$~sebagai !!{sentence}
terakhir dan semua asumsi !!{discharged}. Sebagai contoh, berikut
adalah !!a{derivation} yang menunjukkan bahwa $\Proves (!A
\land !B) \lif !A$:
\begin{prooftree}
  \AxiomC{$\Discharge{!A \land !B}{1}$}
  \RightLabel{\Elim{\land}}
  \UnaryInfC{$!A$}
  \DischargeRule{\Intro{\lif}}{1}
  \UnaryInfC{$(!A \land !B) \lif !A$}
\end{prooftree}
Label~$1$ menunjukkan bahwa asumsi $!A \land !B$ !!{discharged} pada
inferensi \Intro{\lif}.

Himpunan~$\Gamma$ tidak konsisten jika dan hanya jika
$\Gamma \Proves \lfalse$ dalam deduksi alami. Aturan \FalseInt{}
memastikan bahwa setiap !!{sentence} dapat diturunkan dari himpunan
yang tidak konsisten.

Sistem deduksi alami dikembangkan oleh Gerhard Gentzen dan
Stanis\l{}aw Ja\'skowski pada tahun 1930-an, lalu dikembangkan lebih
lanjut oleh Dag Prawitz dan Frederic Fitch. Karena inferensinya
mencerminkan metode pembuktian alami, sistem ini disukai oleh para
filsuf. Versi yang dikembangkan oleh Fitch sering digunakan dalam
buku teks pengantar logika. Dalam filsafat logika, aturan deduksi
alami kadang-kadang dianggap memberikan makna operator logis
(``semantik teori-bukti'').

\end{document}
